% Generated draft: requires mathematical and source review.
\documentclass[11pt]{article}
\usepackage[margin=25mm]{geometry}
\usepackage{fontspec}
\setmainfont{Latin Modern Roman}
\usepackage{amsmath,amssymb,mathtools,unicode-math}
\setmathfont{Latin Modern Math}
\usepackage{enumitem,needspace,xurl,hyperref}
\setlength{\parindent}{0pt}
\setlength{\parskip}{5pt}
\setlength{\emergencystretch}{3em}
\newcommand{\N}{\mathbb{N}}
\newcommand{\Z}{\mathbb{Z}}
\newcommand{\Q}{\mathbb{Q}}
\newcommand{\R}{\mathbb{R}}
\newcommand{\C}{\mathbb{C}}
\newcommand{\F}{\mathbb{F}}
\newcommand{\A}{\mathbb{A}}
\newcommand{\PP}{\mathbb{P}}
\newcommand{\E}{\mathbb{E}}
\newcommand{\eps}{\varepsilon}
\newcommand{\dd}{\,\mathrm{d}}
\begin{document}
\section*{UnsolvedMath ID: 801058}
\textbf{Source problem code:} catalogue entry
\subsection{Definitions and mathematical statement}
Let \(f:\R^2\to\R^2\) be a real H\'enon map. The supplied record does not specify the precise class of maps intended by this term, including the allowed polynomial degrees, normal forms, or compositions; this qualification requires confirmation from the source. No mild dissipativity hypothesis is imposed in the question.

Define the set of periodic points by
\[
\operatorname{Per}(f)=\{x\in\R^2: f^n(x)=x\text{ for some integer }n\geq 1\},
\qquad K_f=\overline{\operatorname{Per}(f)},
\]
where closure is taken in the Euclidean topology of \(\R^2\). The union of all periodic orbits equals \(\operatorname{Per}(f)\). An invariant Borel probability measure is a Borel probability measure \(\mu\) on \(\R^2\) satisfying
\[
\mu(f^{-1}(A))=\mu(A)\qquad\text{for every Borel set }A\subseteq\R^2.
\]
Its support is
\[
\operatorname{supp}\mu=\{x\in\R^2:\mu(U)>0\text{ for every open neighborhood }U\text{ of }x\}.
\]

The question is whether, for every real H\'enon map \(f\) in the intended source class and every invariant Borel probability measure \(\mu\),
\[
\operatorname{supp}\mu\subseteq K_f.
\]
Equivalently, must every such measure satisfy \(\mu(K_f)=1\)? Here ``support all the invariant probability measures'' is interpreted as this support containment, without a compactness assumption on the measures' supports.

As background, the supplied record reports this periodic-point density property for mildly dissipative real H\'enon maps, citing [CP], and for complex H\'enon maps, citing [MR4142445]. These are supplied contextual assertions, not independently verified here. Mild dissipativity is not defined in the record and is not an assumption of the question.
\subsection{Short English statement}
For every real H\'enon map in the source's intended class, must the support of every invariant Borel probability measure lie in the closure of the periodic points? The exact class of real H\'enon maps requires confirmation. The registry labels this question open; this draft does not independently verify that status.
\subsection{Sources}
\begin{itemize}
\item Local UnsolvedMath registry, permanent ID \(801058\), source problem code \texttt{catalogue entry}, title ``Henon Maps: Open Problems --- Question 10.'' The supplied record attributes the problem to Berger and Bedford, with proposed year 2023, and labels its status open.
\item Berger and Bedford, \emph{Henon Maps: A List of Open Problems} (2023), Question 10. Supplied source URL: \url{https://arxiv.org/abs/2312.03907}. The record gives an access date of 2026-10-16 and extraction method \texttt{arxiv\_latex}; these are supplied provenance details, not a claim of source access for this draft.
\item [CP], cited in the supplied statement for the mildly dissipative real case. No full bibliographic entry or link is supplied.
\item [MR4142445], cited in the supplied statement for the complex case. No full bibliographic entry or link is supplied.
\end{itemize}
\end{document}
